\documentclass[11pt]{article}
\usepackage[a4paper,margin=1in]{geometry}
\usepackage{amsmath,amssymb}
\usepackage{booktabs}
\usepackage{graphicx}
\usepackage{enumitem}
\usepackage{fancyhdr}
\usepackage{caption}
\usepackage{longtable}
\usepackage[hidelinks]{hyperref}
\usepackage{xcolor}

\setlength{\headheight}{14pt}
\pagestyle{fancy}
\fancyhf{}
\fancyhead[L]{\small\scshape\nouppercase{\leftmark}}
\fancyhead[R]{\small\itshape\nouppercase{\rightmark}}
\fancyfoot[C]{\thepage}
\renewcommand{\headrulewidth}{0.4pt}

\title{\textbf{\textit{Stock Market Index Direction Prediction Using\\ Limit Order Book Data}}\\[4pt]
\large MTech (DSAI) Capstone --- Final Revision}
\author{Joseph Raj \and Sumanth G.M \and Abhishek Kumar \and Kumara Swamy \and Raghav \and Saagar}
\date{June 2026}

\begin{document}
\maketitle
\thispagestyle{empty}

\begin{abstract}
\noindent We study short-term \emph{index-direction} prediction from Limit Order Book (LOB) data, combining
a faithful reproduction of three published deep LOB models (DeepLOB, MLPLOB, TLOB) on the public FI-2010
benchmark with a first-of-its-kind transfer study on \textbf{NSE index futures (NIFTY / BANKNIFTY)}, for
which we \emph{engineered our own dataset}. We additionally introduce \textbf{MambaLOB}, a selective
state-space (SSM) architecture, and characterise its efficiency--accuracy trade-off against the
transformer baseline. A substantial part of the contribution is \textbf{data engineering}: because no
clean Indian index-futures LOB dataset existed, we reverse-engineered the Dhan binary depth feed, built a
cloud collection pipeline (EC2 $\rightarrow$ S3), and produced a continuous front-month series with a
documented preprocessing pipeline that aligns the NSE data to the FI-2010 contract. This document is the
final interim revision: it reports the reproduction, the engineered NSE benchmark, the architecture
comparison, and an efficiency analysis, with honest framing of what does and does not transfer.
\end{abstract}

\tableofcontents
\clearpage

% =====================================================================
\section{Introduction and Literature Survey}
\markboth{Introduction and Literature Survey}{}

\subsection{Background and Context}
Financial markets facilitate capital allocation, risk distribution, and continuous price discovery.
Early forecasting relied on statistical time-series models applied to closing prices. Modern electronic
exchanges instead operate as continuous double auctions recorded in a \emph{Limit Order Book} (LOB),
which updates at microsecond granularity and produces millions of structured events per session. For an
index such as \textbf{NIFTY 50} --- India's premier equity index --- the LOB of the index-\emph{futures}
contract encapsulates the collective order flow of a broad basket of large-cap stocks, offering a rich
microstructural signal for short-horizon direction prediction.

\subsection{Motivation and Importance of the Study}
Research on deep LOB models is concentrated on developed markets (NASDAQ, Nordic exchanges). Indian index
futures are highly liquid yet structurally distinct (tick discreteness, regime shifts, distinct liquidity
profiles), and remain almost entirely unstudied with modern LOB architectures. Index direction is hard to
predict because the series is non-stationary, exhibits stochastic volatility, and is governed by rapidly
evolving supply--demand dynamics. Crucially, the LOB evolves in \emph{event time}, not clock time: it
updates whenever a limit order is placed, a market order executes, an order is cancelled, or a level is
modified. Most prior deep models resample this irregular stream onto a fixed grid (e.g.\ 1-second
snapshots), which loses intra-interval dynamics, blurs microstructure reactions, and introduces sampling
bias. This study targets Indian index futures directly in event time, addressing a clear geographic and
methodological gap.

\subsection{Problem Definition and Scope}
We predict the direction of the future mid-price over a short horizon from a window of recent LOB states.
The scope comprises: (i) high-frequency LOB inputs (top-$k$ bid/ask prices and volumes); (ii) construction
of an event-aligned NSE index-futures dataset; (iii) reproduction of CNN, MLP, and Transformer baselines;
(iv) a novel selective-SSM architecture (MambaLOB); and (v) multi-horizon, cost-aware evaluation.

\subsection{Research Objectives and Questions}
\textbf{Objectives.}
\begin{itemize}[nosep]
  \item Reproduce DeepLOB and TLOB on FI-2010 and validate against published numbers.
  \item Engineer a reproducible, event-aligned NSE index-futures LOB dataset (NIFTY, BANKNIFTY).
  \item Introduce MambaLOB and characterise its accuracy and computational efficiency.
  \item Compare all architectures on identical data, labels, and protocol; quantify transfer to NSE.
\end{itemize}
\textbf{Research Questions.}
\begin{enumerate}[nosep]
  \item Do our reproductions match published FI-2010 results within tolerance?
  \item Is a linear-time selective SSM (Mamba) competitive with a transformer at lower compute?
  \item Does LOB directional signal transfer from FI-2010 to NSE index futures, and how much does it degrade?
  \item Which labelling scheme (fixed-threshold vs spread-relative) yields an economically meaningful task?
\end{enumerate}

\subsection{Expected Contributions}
\begin{itemize}[nosep]
  \item \textbf{Engineered NSE benchmark.} A continuous front-month NIFTY/BANKNIFTY index-futures LOB
        dataset built from the Dhan depth feed --- to our knowledge the first such Indian index-futures LOB
        dataset prepared for ML, including the full data-engineering pipeline.
  \item \textbf{Reproduction.} Faithful re-implementations of DeepLOB, MLPLOB and TLOB with results
        benchmarked against the original papers and against independent reproductions (LOBCAST).
  \item \textbf{Novel architecture.} MambaLOB, a selective state-space model for LOB, with a rigorous
        efficiency comparison (linear vs quadratic scaling in sequence length).
  \item \textbf{Transfer study.} A standardized comparison of four architectures on the NSE benchmark, with
        honest reporting of degradation relative to FI-2010.
\end{itemize}

\subsection{Review Methodology and Source Selection Criteria}
The survey covers classical statistical approaches (ARIMA), Hidden Markov Models, deep LOB models
(CNN, LSTM, CNN--LSTM hybrids), the transformer-based TLOB/MLPLOB family, state-space models for sequences
(S4/S5/Mamba), and continuous-time microstructure models. Selection criteria: use of high-frequency LOB
data, explicit modelling of LOB structure, empirical validation, and architectural novelty. Primary
references include the DeepLOB family, the TLOB paper, FI-2010 (Ntakaris et al.), the LOBCAST benchmark
study (Prata et al.), the Mamba paper (Gu \& Dao), and continuous-time LOB modelling work.

\subsection{Survey of Existing Approaches}
\subsubsection{Early Solutions and Classical Techniques}
Parametric statistical models (AR, MA, ARMA, ARIMA) assume specific data-generating processes and perform
adequately on linear dynamics but fail under the non-linear, non-stationary behaviour of LOB data. Enhanced
and hybrid HMM/ARIMA models improve regime modelling but depend on historical price series, do not use LOB
depth, and are sensitive to regime misclassification.

\subsubsection{Modern Approaches and Evolving Trends}
\textbf{DeepLOB} (Zhang et al., 2019) combines convolutional feature extraction across LOB levels with an
LSTM over time, substantially improving mid-price movement prediction on FI-2010. It established the
40-feature $\times$ window representation that we adopt. It nonetheless operates on event-resolution
snapshots without attention.

\subsubsection{Emerging Models and Hybrid Frameworks}
\textbf{MLPLOB} and \textbf{TLOB} (Berti \& Kasneci, 2025) reframe the model as alternating normalization and
mixing blocks; TLOB uses \emph{dual} (temporal and feature) self-attention and is the current
state-of-the-art on FI-2010. \textbf{State-space models} (S4, S5, Mamba) offer linear-time sequence
modelling and have entered LOB research on the generative side; their use as a classification backbone for
LOB direction is the novelty window this project targets.

\subsection{Comparative Evaluation of Techniques}
\begin{table}[h]
\centering
\caption{Comparison of LOB modelling approaches considered in this work.}
\begin{tabular}{@{}llll@{}}
\toprule
\textbf{Approach} & \textbf{Strength} & \textbf{Limitation} & \textbf{Role here} \\
\midrule
ARIMA & Linear forecasting & No non-linear/LOB structure & Conceptual baseline \\
HMM & Regime modelling & No microstructure depth & Conceptual baseline \\
DeepLOB (CNN--LSTM) & Depth + temporal & Sequential bottleneck & Reproduced baseline \\
MLPLOB & Lightweight & No sequence model & Ablation floor \\
TLOB (Transformer) & SOTA accuracy & $O(L^2)$ attention cost & Reproduced SOTA baseline \\
MambaLOB (SSM) & $O(L)$, tiny params & Lower accuracy on FI-2010 & \textbf{Proposed model} \\
\bottomrule
\end{tabular}
\end{table}

\subsection{Dataset and Benchmark Overview}
\textbf{FI-2010} (public benchmark): 5 Finnish stocks, 10 trading days, 10-level LOB, pre-normalized
(z-score variant), with standard anchored cross-folds. Used for reproduction.
\textbf{NSE index-futures dataset (engineered here):} NIFTY and BANKNIFTY near-month futures, 20-level depth
captured live via the Dhan API and stored as Parquet on S3; the first ten levels yield a 40-feature vector
identical in layout to FI-2010. Period: 12 May -- 12 June 2026 (21 valid trading days across the May
$\rightarrow$ June expiry roll); $\approx$57k events/instrument/day.

\subsection{Tools, Frameworks, and Experimental Protocols}
Python 3.x; PyTorch ($\geq$2.x) with the CUDA \texttt{mamba-ssm} kernel; scikit-learn (metrics, baselines);
NumPy/Pandas (LOB processing); AWS EC2 + S3 (collection/storage); Google Colab and Kaggle (GPU training);
Matplotlib/Seaborn (figures). Protocol: chronological train/val/test splits; per-feature z-score from
training statistics only; sliding windows of length $L=100$; percentage-mid-change labelling.

\subsection{Evaluation Metrics Used in Literature}
\begin{table}[h]
\centering
\caption{Evaluation metrics for LOB direction prediction.}
\begin{tabular}{@{}lll@{}}
\toprule
\textbf{Metric} & \textbf{Definition} & \textbf{Why it matters} \\
\midrule
Weighted F1 & class-frequency-weighted $F_1$ & matches headline numbers in original papers \\
Macro F1 & unweighted mean $F_1$ & fair under class imbalance; matches independent repro \\
Accuracy & correct/total & supplementary; misleading under imbalance \\
MCC & Matthews correlation & balanced single-number summary \\
Hit-rate / net bps & directional accuracy / cost-adj.\ return & signal-quality / economic relevance \\
\bottomrule
\end{tabular}
\end{table}
We report both weighted and macro $F_1$ throughout, because (as Section~3.6 shows) the original FI-2010
headline numbers are weighted while independent reproductions report macro.

\subsection{Limitations of Prior Work}
No emerging-market (Indian) application; limited event-driven deep architectures; little use of
linear-time state-space models for LOB classification; limited reporting of efficiency/scaling; and
inconsistent metric conventions across papers.

\subsection{Research Gaps and Open Questions}
Can LOB-based models generalise to Indian index futures? Can a linear-time SSM match a transformer at
lower compute? How should labels be defined so that the task is economically meaningful (rather than
trivially dominated by ``no move'')? These motivate our design.

\subsection{Innovation and Conceptual Synthesis}
This work synthesises classical benchmarking, CNN spatial extraction (DeepLOB), transformer dual-attention
(TLOB), and linear-time state-space modelling (Mamba) on a newly engineered Indian benchmark, under one
standardized protocol.

\subsection{Visualizations and Knowledge Summaries}
The report includes comparative tables, architecture descriptions, a data-pipeline schematic, F1-by-horizon
charts, confusion matrices, and the memory/latency-vs-sequence-length scaling curve.

\subsection{Academic Integrity and Citation Style}
Empirical findings are attributed to their source works; conceptual explanations are original synthesis.
Generative-AI assistance (Section~3.11.1) was used for drafting and is disclosed.

\subsection{Summary of the Chapter}
We reviewed classical, deep, transformer, and state-space approaches, identified the Indian-market and
efficiency gaps, and positioned our contribution: an engineered NSE LOB benchmark, a reproduction of three
baselines, and a novel Mamba-based model evaluated on identical footing.

\subsection{Structure of the Report}
Chapter~1 (this chapter) introduces the problem and surveys the literature. Chapter~2 establishes the
theoretical and architectural foundations. Chapter~3 is the interim report: the engineered dataset, the
data-engineering pipeline, preprocessing, baseline development, preliminary results, efficiency analysis,
reproducibility, risks, and next steps.

% =====================================================================
\section{Foundations}
\markboth{Foundations}{}

\subsection{Introduction to the Foundations Chapter}
This chapter establishes the theoretical underpinnings (market microstructure, the limits of market
efficiency), the conceptual data-to-model pipeline, the four model architectures (three reproduced
baselines and the proposed MambaLOB), and the notation used throughout.

\subsection{Theoretical Underpinnings}
\subsubsection{Market Microstructure}
The continuous double auction maintains a LOB recording all outstanding limit orders at each price level.
The \emph{bid} side (buy orders, descending) has best bid $P_1^b$; the \emph{ask} side (sell orders,
ascending) has best ask $P_1^a$. The mid-price is the primary, less noisy, forecasting target:
\begin{equation}
p_{\mathrm{mid}}(t) = \frac{P_1^a(t) + P_1^b(t)}{2}.
\end{equation}
The bid--ask spread $S(t) = P_1^a(t) - P_1^b(t)$ measures liquidity (and, importantly for us, transaction
cost), while per-level depth encodes short-term price pressure.

\subsubsection{The Efficient Market Hypothesis and Its Limits}
The semi-strong EMH implies no excess return from public information, yet high-frequency studies document
short-lived predictability from order-flow imbalance, inventory effects, and latency arbitrage at
sub-second to minute horizons. Deep models trained on LOB snapshots exploit precisely these transient
inefficiencies; their documented (if small) predictive power motivates this study, while the EMH frames
our expectation that signal is weak and easily overwhelmed by costs.

\subsection{Conceptual Framework}
The pipeline is: \emph{collect} (FI-2010 download; live Dhan capture for NSE) $\rightarrow$ \emph{preprocess}
(clean, reorder to a 40-feature LOB vector, z-score on train statistics, label, window) $\rightarrow$
\emph{model} (DeepLOB, MLPLOB, TLOB baselines; MambaLOB proposed) $\rightarrow$ \emph{evaluate}
(multi-horizon classification metrics, efficiency, cost-aware backtest). Every model consumes the same
tensor $X \in \mathbb{R}^{B \times L \times F}$ ($L=100$, $F=40$) and emits logits over three classes.

\subsubsection{Feature Representation}
Each LOB state is the first ten levels as
$[\,P^a_i, V^a_i, P^b_i, V^b_i\,]_{i=1}^{10}$ (40 features), the FI-2010 convention. Mid-price and spread
are derived from level~1 for labelling and cost accounting.

\subsubsection{Baseline Models}
\textbf{DeepLOB} (CNN + Inception + LSTM); \textbf{MLPLOB} (flatten + MLP + Bilinear Normalization, the
ablation floor); \textbf{TLOB} (dual temporal/feature self-attention, SOTA). All are reproduced faithfully.

\subsubsection{Proposed MambaLOB Model}
MambaLOB projects the 40-feature input to $d_{\mathrm{model}}$, applies a stack of selective state-space
(Mamba) blocks for $O(L)$ temporal mixing, optionally a feature-axis attention, then Bilinear Normalization
and a linear classifier:
\[
X \xrightarrow{\text{Linear}} \text{Mamba}\times n \xrightarrow{(\text{opt.\ feature attn})}
\text{BiN} \xrightarrow{\text{FC}} \text{logits}\in\mathbb{R}^{3}.
\]
On GPU it uses the CUDA \texttt{mamba-ssm} kernel; a pure-PyTorch fallback exists for CPU/debug.

\subsubsection{Training and Evaluation}
Supervised three-class classification; class-weighted cross-entropy; AdamW with cosine decay; early
stopping on validation macro/weighted-$F_1$; metrics as in Table~2.

\subsection{Terminology and Notational Conventions}
\begin{table}[h]
\centering
\caption{Notation.}
\begin{tabular}{@{}ll@{}}
\toprule
Symbol & Meaning \\
\midrule
$X \in \mathbb{R}^{B\times L\times F}$ & input window batch ($L=100$, $F=40$) \\
$L, F, B$ & sequence length, features, batch size \\
$d_{\mathrm{model}}, d_{\mathrm{state}}, d_{\mathrm{conv}}$ & Mamba projection / state / conv width \\
$k$ (or $H$) & prediction horizon in LOB events \\
$p_{\mathrm{mid}}, S$ & mid-price, bid--ask spread \\
$\theta, \alpha$ & label thresholds (spread-relative / fixed) \\
\bottomrule
\end{tabular}
\end{table}

\subsubsection{Class Labels}
We use the encoding $\{0=\text{Down},\,1=\text{Stationary},\,2=\text{Up}\}$. Labels are assigned by the
smoothed-mid percentage change over the horizon (Section~2.6).

\subsection{Algorithmic Foundations}
\subsubsection{Data Preprocessing Algorithms}
\textbf{Z-score normalization} per feature, using \emph{training} statistics only:
$Z = (X-\mu)/\sigma$.
\textbf{Sliding-window} segmentation produces $X\in\mathbb{R}^{L\times F}$ with $L=100$.
\textbf{Labelling}: with $m^-(t)$ and $m^+(t)$ the mean mid over the $k$ events behind/ahead,
$r(t) = (m^+ - m^-)/m^-$; then Up if $r>\tau$, Down if $r<-\tau$, else Stationary, where $\tau$ is the
threshold (Scheme A: fixed $\alpha$; Scheme B: spread-relative $\theta$).

\subsubsection{DeepLOB / MLPLOB / TLOB}
DeepLOB stacks $1\times2$ convolutions (fuse price/volume, then sides), a $1\times10$ convolution (fuse
levels), an Inception module, and an LSTM. MLPLOB flattens and applies MLP blocks with Bilinear
Normalization. TLOB alternates temporal and feature self-attention with a progressive MLP head.

\subsubsection{Bilinear Normalization and Attention}
Bilinear Normalization (BiN) normalises across both the temporal and feature axes and learns to weight the
two, stabilising LOB training. Multi-head self-attention provides cross-position mixing in TLOB; its cost
is quadratic in $L$.

\subsubsection{Selective State-Space Modelling}
A Mamba block realises a discretised, input-dependent linear recurrence
$h_t = \bar{A}_t h_{t-1} + \bar{B}_t x_t,\ y_t = C_t h_t$, where the selection mechanism makes
$\bar A,\bar B,C$ depend on the input. This yields $O(L)$ time and memory and a parameter count independent
of $L$ --- the basis of our efficiency claim.

\subsection{Classification Algorithm}
The head produces $\hat{y} = \mathrm{Softmax}(Wx+b)$ over the three classes. Optimization uses
class-weighted cross-entropy and AdamW; evaluation uses weighted/macro $F_1$, accuracy, and MCC.

\subsection{Model Design Principles}
Modularity (separate ingest / dataset / model / train / eval modules), reusability (shared loaders,
normalization, windowing across FI-2010 and NSE), adaptability (one data contract for all models), and
scalability (lazy windowing keeps memory $O(T\cdot F)$ rather than materialising all windows).

\subsection{Domain-specific Constraints and Assumptions}
Tick discreteness (NIFTY futures move in $\geq$0.05 increments) quantizes mid-price changes and informs
threshold choice; the bid--ask spread sets a hard floor on tradeable signal; pre-open and the
opening/closing auction windows are excluded to isolate continuous trading; expiry days require contract
rolling.

\subsection{Ethical, Legal, and Societal Implications}
The work uses public benchmark data and self-collected public market data (no personal data). Predictions
are framed as \emph{signal-quality} research, not trading advice; the cost-aware analysis (Section~3.10)
explicitly guards against over-claiming profitability.

\subsection{Summary of the Chapter}
We established microstructure foundations, the shared data-to-model pipeline, the four architectures, and
the notation. The next chapter reports the engineered dataset and the empirical results.

% =====================================================================
\section{Interim Report}
\markboth{Interim Report}{}

\subsection{Chapter Overview}
This chapter documents the work completed to date: (i) the \textbf{data-engineering} effort to build an NSE
index-futures LOB dataset; (ii) the preprocessing pipeline aligning it to FI-2010; (iii) reproduction of
the three baselines and the proposed MambaLOB; (iv) preliminary results on FI-2010 and the NSE transfer;
(v) an efficiency analysis; and (vi) reproducibility, risks, and next steps. Experiments that are still
running or deferred are marked accordingly.

\subsection{Dataset Acquisition and Preprocessing}

\subsubsection{Dataset Description}
\textbf{FI-2010 (reproduction).} The public NoAuction z-score variant: 5 stocks, 10 days, 10 levels,
pre-normalized, with standard cross-folds (we use CF\_7: 7 train + 3 test days).

\textbf{NSE index futures (engineered).} We built a live dataset for NIFTY and BANKNIFTY near-month futures
using the Dhan twenty-level depth feed, stored as daily Parquet on Amazon S3. Key facts: 21 valid trading
days (12 May -- 12 June 2026), spanning the May$\rightarrow$June expiry roll; $\approx$57{,}000
events/instrument/day ($\approx$1.1M combined per instrument, exceeding FI-2010's $\approx$400k scale); the
first 10 of 20 levels give a 40-feature vector identical in layout to FI-2010, so models transfer with no
architectural change.

\subsubsection{Data Quality Checks}
We evaluated two Indian data sources before committing. Table~\ref{tab:dhanvskite} summarises the
comparison that led us to choose Dhan.
\begin{table}[h]
\centering
\caption{Dhan vs Kite vs FI-2010 (15-min simultaneous capture).}
\label{tab:dhanvskite}
\begin{tabular}{@{}llll@{}}
\toprule
Property & Dhan (20-depth) & Kite Connect & FI-2010 \\
\midrule
LOB levels & 20 & 5 & 10 \\
Features (10-level) & 40 (direct match) & 20 (needs change) & 40 \\
Update mechanism & event-driven (gap CV 3.3) & periodic $\sim$1s (CV 1.3) & event-driven \\
Median inter-tick gap & 0.20 s & 1.0 s & $<$0.1 s \\
L1--L10 completeness & 100\% & n/a & 100\% \\
Order-count field & yes & no & no \\
Cost & low & higher & free \\
\bottomrule
\end{tabular}
\end{table}
Dhan wins on depth, event-driven granularity (structurally consistent with FI-2010), completeness, and
cost. Beyond source selection, per-day cleaning drops: rows with crossed quotes ($P^b_1\!\ge\!P^a_1$);
rows with any missing level-1--10 field; and tick-level glitches (mid-price jumps $>$0.5\% relative to the
previous valid mid). Empty/partial collection days (the broken expiry-roll files of 27--29 May) are
excluded; the partial 11 May session is dropped.

\subsubsection{Exploratory Data Analysis (EDA)}
Active instruments reach $\approx$200--250 events/min in single-instrument capture; the heavy-tailed
inter-tick gap distribution (coefficient of variation $\approx$3.3) confirms genuine \emph{event-driven}
updates rather than periodic polling. Median spreads are a few basis points (NIFTY $\approx$2.5 bps),
which --- as Section~3.6 shows --- has direct consequences for labelling. Class balance varies sharply with
horizon and labelling scheme (Tables in Section~3.6).

\subsubsection{Pre-processing Pipelines}
The pipeline that aligns Dhan to the FI-2010 contract (implemented in \texttt{modeling/nse\_dataset.py}):
\begin{enumerate}[nosep]
  \item \textbf{Front-month stitching}: match each daily file by root (NIFTY/BANKNIFTY), concatenating the
        May and June contracts into one continuous front-month series.
  \item \textbf{Feature reorder}: map Dhan's column ranges to the interleaved FI-2010 layout
        $[P^a_i,V^a_i,P^b_i,V^b_i]\times 10$.
  \item \textbf{Cleaning}: crossed-quote, missing-level, and tick-spike filters (above).
  \item \textbf{Per-day segmentation}: labels and windows are computed \emph{within} a trading day, so the
        overnight gap and the expiry roll are segment boundaries --- no spurious cross-day returns and no
        window straddles two contracts.
  \item \textbf{Z-score} using training statistics only.
  \item \textbf{Labelling} (Scheme A fixed-$\alpha$ or Scheme B spread-relative).
  \item \textbf{Lazy sliding windows}: the $(T,F)$ matrix is stored once and windows are sliced on demand,
        keeping memory $O(T\cdot F)$; the eager $(N,L,F)$ expansion would exhaust memory on the
        full series.
\end{enumerate}
Default split (chronological): train 12 May -- 4 June (15 days), validation = last 10\% of training
windows, test 5/8/9/10/11/12 June (6 most-recent out-of-sample days).

\subsection{Feature Engineering}
\subsubsection{Feature Identification Techniques}
The model input is the raw 40-dimensional LOB vector (10 levels) per event; mid-price and spread are
derived for labelling and cost accounting. We deliberately avoid hand-crafted trading indicators, letting
the sequence model learn from the order-book windows directly (consistent with DeepLOB/TLOB).

\subsubsection{Dimensionality Reduction}
None applied: the input is intentionally the raw LOB tensor. We instead study the \emph{depth} choice
(first 10 of Dhan's 20 levels, to match FI-2010) and note Dhan's additional order-count field and levels
11--20 as available future feature channels (an 80-feature or 60-feature variant).

\subsubsection{Feature Impact Observations}
The most consequential ``feature engineering'' decision is the labelling threshold (Section~3.6): under a
spread-relative threshold most NSE moves do not clear the spread, which is itself an economic finding about
signal availability.

\subsection{Technical Implementation}
\subsubsection{Tools and Technology Stack}
Python; PyTorch with \texttt{mamba-ssm}/\texttt{causal-conv1d} (CUDA); scikit-learn; NumPy/Pandas; Dhan and
Kite WebSocket APIs; AWS EC2 (collection) and S3 (storage, region \texttt{ap-south-2}); Google Colab and
Kaggle (GPU training); \texttt{boto3}/\texttt{s3fs}; Matplotlib/Seaborn; Git/GitHub for version control.

\subsubsection{System Setup and Environment Configuration (Data Engineering)}
Because no clean Indian index-futures LOB dataset existed, the dataset was \emph{engineered end to end}:
\begin{itemize}[nosep]
  \item \textbf{Reverse-engineering the Dhan binary feed}: the depth WebSocket emits an undocumented binary
        packet. We decoded its structure offline --- a 1992-byte full format (40-byte header; 20$\times$16-byte
        bid then ask records of order-count/price(f64)/qty) and a 664-byte compact format of two 332-byte
        sub-packets --- fixing byte-offset bugs (bid price at $+4$, not $+0$) and rejecting denormalized
        near-zero floats from empty levels.
  \item \textbf{Collection pipeline}: an EC2 (Ubuntu) collector streams 20-level depth during market hours
        (09:15--15:30 IST), writes daily zstd Parquet, and a cron job syncs to S3 after close.
  \item \textbf{Production debugging}: corrected the futures segment code (\texttt{NSE\_FNO}); diagnosed a
        mislabelled-instrument bug caused by a broker token reassignment after a 1:1 bonus issue; handled
        graceful WebSocket close at market end.
  \item \textbf{Contract roll}: the May contract expired and collection rolled to the June contract; the
        loader stitches them into a continuous front-month series with the roll as a segment boundary.
  \item \textbf{Training environments}: FI-2010 reproduction and Mamba/NSE training run on Colab/Kaggle
        GPUs; secrets (GitHub PAT, AWS keys, Dhan/Kite tokens) are read host-agnostically; every experiment
        uploads its results and checkpoints to S3 as it completes (timeout-safe, resumable across sessions).
\end{itemize}

\subsubsection{Modular Code-base Description}
The repository (\texttt{nse-lob-capstone}) is organised as: \texttt{data\_acquisition/} (Dhan/Kite
collectors, token refresh, S3 sync, EC2 ops); \texttt{modeling/} (\texttt{fi2010\_dataset.py},
\texttt{nse\_dataset.py}, \texttt{models.py} [DeepLOB/MLPLOB/MambaLOB], \texttt{tlob.py},
\texttt{train.py}, \texttt{backtest.py}, \texttt{stats.py}, \texttt{nbenv.py}); \texttt{notebooks/}
(reproduction, TLOB, Mamba+efficiency, NSE matrix, NSE extras); \texttt{literature/}; and \texttt{docs/}.
All models consume one data contract, so a single loader/trainer serves both datasets.

\subsection{Baseline Model Development}
\subsubsection{Baseline Algorithm Description}
DeepLOB (CNN--Inception--LSTM), MLPLOB (MLP+BiN), and TLOB (dual-attention transformer) are implemented to
the shared $(B,100,40)\rightarrow[B,3]$ contract; TLOB is vendored faithfully from the official
implementation.

\subsubsection{Training Setup and Parameters}
Max 50 epochs (FI-2010) / 20 epochs (NSE), early stopping on validation $F_1$; AdamW; class-weighted
cross-entropy; batch 128; per-model learning rates following the papers (DeepLOB/MLPLOB $10^{-3}$, TLOB
$10^{-4}$, MambaLOB $3\times10^{-4}$).

\subsubsection{Early Performance Metrics --- FI-2010 reproduction}
\begin{table}[h]
\centering
\caption{FI-2010 (CF\_7, $L=100$) --- weighted / macro $F_1$ (\%).}
\label{tab:fi2010}
\begin{tabular}{@{}lccccc@{}}
\toprule
Model & $k{=}10$ & $k{=}20$ & $k{=}50$ & $k{=}100$ & Notes \\
\midrule
DeepLOB & 82.5 / 71.9 & --- & --- & --- & $\approx$ paper 83.4 (weighted) \\
TLOB    & 88.0 / 81.1 & 77.3 / 70.2 & 88.4 / 87.5 & 92.5 / 92.5 & SOTA-range; paper $\sim$92.8 @ seq 128 \\
MambaLOB (M1) & 79.2 / 67.4 & 71.1 / 61.4 & 72.8 / 70.0 & 69.3 / 68.7 & lower accuracy, far cheaper \\
\bottomrule
\end{tabular}
\end{table}
\textbf{Acceptance.} DeepLOB's weighted $F_1$ (82.5\%) matches the paper's headline ($\sim$83.4\%); TLOB is
in the SOTA range, reaching 92.5\% at $k{=}100$ (the paper's $\sim$92.8 was obtained at sequence length
128; we use 100). These confirm the reproduction is correct.

\subsubsection{Output Visualization and Interpretation}
We report confusion matrices per (model, horizon), $F_1$-by-horizon line charts, and the efficiency
scaling curve (Section~3.10). At longer horizons macro and weighted $F_1$ converge as the Stationary class
shrinks.

\subsection{Preliminary Analysis and Observations}
\subsubsection{Result Analysis and Discussion}
\textbf{FI-2010.} TLOB $>$ DeepLOB $>$ MambaLOB on accuracy. \textbf{NSE transfer (Scheme A, partial).}
On NIFTY, weighted $F_1$ is $\approx$0.44--0.52 (DeepLOB), 0.46--0.64 (MLPLOB) and 0.60--0.69 (TLOB); all
models beat the naive baselines (random $\approx$0.33), so there \emph{is} learnable signal --- but the
absolute numbers are well below FI-2010, the expected degradation on new data (consistent with the LOBCAST
benchmark study). MambaLOB's NSE rows are still training at the time of writing.

\subsubsection{Error Profiling and Bias Detection}
The dominant effect is class balance. Under the fixed threshold the Stationary class is small at long
horizons (near-binary), whereas under the spread-relative scheme it dominates. Confusion matrices show
errors concentrated on the minority directional classes, which is why macro $F_1$ trails weighted $F_1$.

\subsubsection{Lessons Learned}
(1) \textbf{Metric convention matters}: the FI-2010 headline ($\sim$83.4) is a \emph{weighted} $F_1$; our
weighted $F_1$ matches it, while our macro $F_1$ ($\approx$72) matches \emph{independent} reproductions
(LOBCAST). We therefore report both. (2) \textbf{F1 $\ne$ profit}: see Section~3.6 / 3.10. (3)
\textbf{Mamba is an efficiency win, not an accuracy win} on FI-2010.

\subsection{Originality and Innovation in Approach}
\subsubsection{Creative Framing of the Problem}
The first ML-ready NSE index-futures LOB benchmark, framed honestly for snapshot depth data (an ``event''
is a change in the top-of-book), and a transaction-cost-relative labelling scheme that exposes when signal
is economically meaningful.
\subsubsection{Experimental Innovations}
MambaLOB (selective SSM for LOB classification); a controlled long-context experiment exploiting Mamba's
linear scaling; and two labelling schemes (fixed vs spread-relative) reported side by side.

\subsection{Risk Assessment and Project Management}
\subsubsection{Progress Tracker vs.\ Original Timeline}
Completed: literature survey; the engineered NSE dataset + pipeline; FI-2010 reproduction (DeepLOB, MLPLOB,
TLOB); MambaLOB on FI-2010 + efficiency analysis; NSE Scheme-A matrix (DeepLOB/MLPLOB/TLOB done, Mamba in
progress). Deferred (next phase): Scheme-B confirmatory training, multi-seed confidence intervals, and the
cost-aware backtest.
\subsubsection{Bottlenecks and Delays}
Free-tier GPU throttling (mitigated by moving from Colab to Kaggle, which has a separate quota); a Kaggle
P100 incompatible with the latest PyTorch CUDA build (use T4); \texttt{mamba-ssm} failing to build under
pip build isolation; and session time-outs on long matrices.
\subsubsection{Risk Mitigation Strategies}
\texttt{mamba-ssm} installed with \texttt{--no-build-isolation}; per-run S3 upload + pull-on-start makes
every long matrix resumable (a time-out costs at most the run in flight); host-agnostic secrets/data so the
same notebooks run on Colab or Kaggle.

\subsection{Experimental Reproducibility}
\subsubsection{Data Splits and Control Measures}
FI-2010: standard CF\_7 (7 train / 1 val / 3 test days). NSE: chronological front-month split (train
12 May--4 June; test 5--12 June). Normalization statistics are computed on training data only; windows and
labels never cross a day/contract boundary.
\subsubsection{Random Seed Control}
A \texttt{set\_seed} utility seeds Python/NumPy/PyTorch; the headline configurations are scheduled for
3-seed runs with bootstrap 95\% confidence intervals (deferred to the next phase). All results, metrics
JSON, and checkpoints are versioned on S3 under \texttt{reproduction/}.

\subsection{Scalability and Efficiency Considerations}
\subsubsection{Memory, Time, and Resource Profiling}
\begin{table}[h]
\centering
\caption{Efficiency at $L{=}100$, batch 1 (inference).}
\begin{tabular}{@{}lccc@{}}
\toprule
Model & Params & Latency (ms) & Peak mem (MB) \\
\midrule
DeepLOB & 191k & 1.5 & 30 \\
MLPLOB & 529k & 0.2 & 22 \\
TLOB & 1.80M & 6.1 & 27 \\
MambaLOB & \textbf{68k} & 1.5 & 20 \\
\bottomrule
\end{tabular}
\end{table}
\begin{table}[h]
\centering
\caption{Scaling with sequence length (batch 32): TLOB $O(L^2)$ vs MambaLOB $O(L)$.}
\begin{tabular}{@{}lcccc@{}}
\toprule
$L$ & TLOB params & TLOB mem / latency & Mamba params & Mamba mem / latency \\
\midrule
100 & 1.8M  & 49 MB / 11 ms  & 68k & 32 MB / 1.5 ms \\
400 & 13.5M & 164 MB / 54 ms & 68k & 69 MB / 1.8 ms \\
800 & 51M   & 400 MB / 164 ms & 68k & 118 MB / 5.3 ms \\
\bottomrule
\end{tabular}
\end{table}
\textbf{Headline efficiency result.} MambaLOB's parameter count is constant in $L$ and its
memory/latency grow \emph{linearly}, while TLOB grows \emph{quadratically}. At $L{=}800$ MambaLOB is
$\approx$31$\times$ faster, uses $\approx$3.4$\times$ less memory, and has $\approx$750$\times$ fewer
parameters --- a clean demonstration of the linear-scaling advantage that motivates the architecture.
(FLOP counts for Mamba are under-reported by the profiler because the fused CUDA kernel is not traced; we
therefore rely on measured latency and memory.)
\subsubsection{Scalability Constraints}
Training throughput is bounded by data loading and the sequential scan rather than raw FLOPs; the
spread-relative task is degenerate at short horizons; and real-time deployment is out of scope.

\subsection{Ethical Considerations and Responsible AI Use}
\subsubsection{Generative AI Usage Disclosure}
Generative-AI assistance (a coding/writing assistant) was used for code scaffolding, debugging, drafting
documentation, and report formatting. All generated content and code were reviewed, tested, and validated
before inclusion.
\subsubsection{Bias and Fairness in Dataset}
The data are public benchmark data and self-collected public market data; there is no personal or
demographic data. The main ``bias'' risk is class imbalance, which we address with class-weighted loss and
by reporting macro $F_1$ and baselines.

\subsection{Feedback Incorporation and Continuous Improvement}
\subsubsection{Peer and Mentor Feedback Summary}
Feedback drove: separating weighted vs macro $F_1$ reporting; checking the spread-relative label balance
before training (which revealed degeneracy and saved a large, pointless run); and prioritising the
data-engineering narrative.
\subsubsection{Reflection and Learning Journey}
The project built skills across data engineering (binary protocol decoding, cloud collection, robust
preprocessing), reproducible deep-learning experimentation, and honest empirical evaluation under cost and
compute constraints.

\subsection{Chapter Summary and Next Steps}
We engineered an NSE index-futures LOB dataset, reproduced three baselines on FI-2010, introduced MambaLOB
with a strong efficiency result, and began the NSE transfer study. \textbf{Next steps}: finish the
MambaLOB NSE runs; the Scheme-B confirmatory training; multi-seed bootstrap confidence intervals; the
cost-aware backtest; a walk-forward robustness check; and consolidation of all results into the final
report.

\paragraph{Appendices (planned).} Detailed per-(model, symbol, horizon) result tables; confusion-matrix and
scaling figures; the notebook$\rightarrow$result map; and reproducibility artefacts (\texttt{requirements},
seeds, S3 paths).

\end{document}
